% Generated draft: requires mathematical and source review.
\documentclass[11pt]{article}
\usepackage[margin=25mm]{geometry}
\usepackage{fontspec}
\setmainfont{Latin Modern Roman}
\usepackage{amsmath,amssymb,mathtools,unicode-math}
\setmathfont{Latin Modern Math}
\usepackage{enumitem,needspace,xurl,hyperref}
\setlength{\parindent}{0pt}
\setlength{\parskip}{5pt}
\setlength{\emergencystretch}{3em}
\newcommand{\N}{\mathbb{N}}
\newcommand{\Z}{\mathbb{Z}}
\newcommand{\Q}{\mathbb{Q}}
\newcommand{\R}{\mathbb{R}}
\newcommand{\C}{\mathbb{C}}
\newcommand{\F}{\mathbb{F}}
\newcommand{\A}{\mathbb{A}}
\newcommand{\PP}{\mathbb{P}}
\newcommand{\E}{\mathbb{E}}
\newcommand{\eps}{\varepsilon}
\newcommand{\dd}{\,\mathrm{d}}
\begin{document}
\section*{UnsolvedMath ID: 20002042}
\textbf{Source problem code:} catalogue entry
\subsection{Definitions and mathematical statement}
This merged extraction contains two questions. Its first piece is a tail referring to Problem 4 in dimension six and a proposed decomposition
\[
 L=c_L\operatorname{Pff}+\operatorname{div}_L+I_L,
\]
where \(I_L\) is a local conformal invariant and the terms depend on \(L\). It attributes to T. Branson the observation that the coefficient for \(Q\) is \(c_Q\ne0\), computable using the constant sphere values of \(Q\) and the Pfaffian. The aid explicitly corrects the OCR fragment ``constant Q 6 =0'' to nonzero, not to a six-dimensional vanishing assertion. The preceding Problem 4, the class of allowed natural densities, and hypotheses such as conformal invariance of the integral are absent. Density weight \(-n\) alone is not stated here as a sufficient hypothesis for the standard decomposition theorem.

The second piece asks how \(Q\)-curvature relates to Weyl structures. A Weyl structure is a conformal class with a torsion-free compatible connection. In one explicit sign convention, \(Dg=2\alpha\otimes g\), so for \(\widehat g=e^{2\varphi}g\), \(\widehat\alpha=\alpha+\dd\varphi\). The Faraday two-form \(F=\dd\alpha\) is independent of representative, and the Weyl structure is closed when \(F=0\); closed does not require a globally exact \(\alpha\).

The source treats \(Q\) as a top-degree density and factors the critical operator as
\[
 P f=S_g(\dd f),
\]
where \(S_g\) is a natural linear operator from one-forms to top forms. It considers \(Q_g-S_g\alpha\) and records a conformal-change formula that leaves a dependence on \(S_{\widehat g}\); thus this expression is not simply asserted to be invariant without correction. In dimension four it attributes to Eastwood--Slovák the corrected expression, in the source's normalization,
\[
 Q_{\mathrm{ES}}=Q_g-S_g\alpha
 +4\nabla^a\bigl(\alpha^b\nabla_{[a}\alpha_{b]}\bigr)\dd V_g.
\]
The source says it is an invariant of the Weyl structure agreeing with ordinary \(Q\) for a Riemannian Weyl structure. The exact \(Q,P,S\), bracket, and density conventions are not fully reproduced, so this formula is retained with that normalization qualification.

The aid attributes the standard natural polynomial decomposition to Alexakis and a Weyl \(Q\)-curvature cotractor construction in every even dimension \(n\ge4\) to Hirachi--Lubbe--Matsumoto. It supplies no full theorem citations or cotractor formula. Its comparison theorem is in dimension four: let \(M\) be a compact oriented Riemannian four-manifold without boundary and normalize the Eastwood--Slovák and Hirachi--Lubbe--Matsumoto extensions to the same Branson \(Q\)-curvature. In the report's common conventions it gives
\[
 \int_M Q_{\mathrm{ES}}=\int_M Q_g,\qquad
 \int_M Q_{\mathrm{HLM}}=\int_M Q_g+2\int_M|F|_g^2\dd V_g,
\]
and hence
\[
 \int_M(Q_{\mathrm{HLM}}-Q_{\mathrm{ES}})
 =2\|F\|_{L^2(g)}^2.
\]
The numerical factor is convention dependent and is quoted in the aid's normalization, not asserted across other \(Q\) or two-form norm conventions. For \(F\not\equiv0\), the difference cannot be a pure divergence on such a manifold. The aid also reports agreement for closed Weyl structures; zero integrated discrepancy alone would not establish pointwise equality without the construction-level comparison.

The aid stresses that in the HLM argument only the operator named \(L_1\) is asserted to annihilate closed forms. The \(G_1\) contribution vanishes after integration because it is a divergence; it must not be said to annihilate closed forms pointwise. The operators are not defined in the supplied record. The resulting conclusion is nonuniqueness of natural Weyl \(Q\)-extensions away from closed structures, not failure of the cited construction or a resolution of every unqualified interpretation of the omitted Pfaffian problem.
\subsection{Short English statement}
The record combines a Pfaffian-decomposition fragment with a question about extending \(Q\)-curvature to Weyl connections. The aid reports two natural extensions with different integrals when the Faraday form is nonzero, establishing nonuniqueness in its stated conventions. It preserves the nonzero OCR correction and separates an operator's pointwise property from a divergence that vanishes only after integration.
\subsection{Sources}
\begin{itemize}
\item Local UnsolvedMath registry, permanent ID 20002042, exact source code catalogue entry; canonical \texttt{aim-geometry-notes.json}, \texttt{notes[379]}, tag \texttt{problem}; AIM Workshop Problem Lists, Geometry. Created and updated 14 August 2026.
\item AIM 2003 workshop document \emph{Conformal Structure in Geometry, Analysis, and Physics}; \url{https://aimath.org/WWN/confstruct/confstruct.pdf}. The supplied verification identifies two logically separate extracted pieces. Original attributions: T. Branson for the nonzero Pfaffian coefficient, Eastwood and Slovák for the four-dimensional Weyl formula. The PDF was not accessed.
\item The aid names Alexakis for the standard decomposition and Hirachi--Lubbe--Matsumoto for the even-dimensional Weyl \(Q\)-cotractor. No titles, publication identifiers, theorem numbers, or separate URLs are supplied.
\item Attempt 1: \texttt{partial\_result}, problem status at run \texttt{solved}, rigor label \texttt{proved}, novelty confidence \texttt{low}; classification \texttt{SOLVED-IN-LITERATURE}, registry status \texttt{solved}. Candidate: four-dimensional integrated comparison and nonuniqueness. These are dated source assessments, not independent certification.
\item Ulam AI / UnsolvedMath Contributors; CC BY 4.0. Machine-generated research aid, not peer reviewed; independently verify before citation.
\end{itemize}
\end{document}
